\documentclass[10pt]{article}
\providecommand{\FIGDIR}{fig}
\newcommand{\DIGESTDATE}{19 October 2026}
\input{preamble}

\begin{document}

% ==================================================================== MASTHEAD
\thispagestyle{empty}
\begin{tikzpicture}[remember picture, overlay]
  \fill[Deep] (current page.north west) rectangle
      ($(current page.north east)+(0,-67mm)$);
  \fill[Amber] ($(current page.north west)+(0,-67mm)$) rectangle
      ($(current page.north east)+(0,-68.4mm)$);
  \foreach \x/\y/\r in {22/12/0.9, 41/26/1.5, 63/9/0.7, 88/31/1.1, 112/17/1.9,
                        138/29/0.8, 159/13/1.3, 178/54/1.0, 196/21/1.6,
                        30/44/0.6, 74/57/1.2, 125/48/0.9, 168/6/0.8, 191/59/0.7,
                        52/18/1.0, 100/60/0.6, 148/20/1.1, 66/49/1.4, 205/45/0.8}
    \fill[Sky, opacity=0.30] ($(current page.north west)+(\x mm,-\y mm)$) circle (\r mm);
\end{tikzpicture}

\vspace*{4mm}
{\sffamily\fontsize{7.6}{9}\selectfont\color{Sky}%
 AUTOMATED DAILY LITERATURE DIGEST\;\textbullet\;PhD RESEARCH WATCH\par}
\vspace{2.2mm}
{\sffamily\bfseries\fontsize{27}{29}\selectfont\color{white}%
 Particulate Matter\par}
{\sffamily\fontsize{27}{29}\selectfont\color{Sky}%
 Monitoring \& Health Effects\par}

\vspace{5.0mm}
{\sffamily\fontsize{8.4}{10}\selectfont\color{white}%
 Issue 2026-10-19\;\;\textcolor{Amber}{\textbar}\;\;Entry window 19 October 2026 (full day)\;\;%
 \textcolor{Amber}{\textbar}\;\;18 records screened in scope\par}
\vspace{18mm}

\noindent
\begin{minipage}[t]{0.190\linewidth}\metric{18}{RECORDS IN SCOPE\\(70 EXCLUDED)}{Deep}\end{minipage}\hfill
\begin{minipage}[t]{0.190\linewidth}\metric{9}{OF 10 SUBTOPIC\\BANDS POPULATED}{Teal}\end{minipage}\hfill
\begin{minipage}[t]{0.190\linewidth}\metric{5}{EFFECT ESTIMATES\\WITH A CI}{Sky}\end{minipage}\hfill
\begin{minipage}[t]{0.190\linewidth}\metric{1}{TIER-A\\RECORDS}{Amber}\end{minipage}\hfill
\begin{minipage}[t]{0.190\linewidth}\metric{6}{RECORDS CARRIED ON\\METADATA ALONE}{Clay}\end{minipage}

\vspace{6mm}

\section*{\sffamily\bfseries\color{Deep}Signal of the day}
\vspace{-1.5mm}{\color{Amber}\rule{\linewidth}{1.1pt}}\vspace{2.5mm}

\begin{keybox}
{\sffamily\bfseries\small\color{Deep}In 131,346 Guangdong schoolchildren, mass-normalised oxidative potential
explains allergy-related symptoms more consistently than PM$_{2.5}$ mass --- and carries up to 36\,\% of a
constituent's association.}\\[1.4mm]
\footnotesize
Six cities, ages 6--18; DTT oxidative potential as volume- (DTTv) and mass-normalised (DTTm) metrics alongside
PM$_{2.5}$ mass, metals, carbonaceous fractions and ions. Mass and OP were both positively associated with rhinitis,
rhinoconjunctivitis and eczema symptoms; DTTm gave lower AIC/BIC and survived two-pollutant models best. Mixture
and mediation analyses put the largest indirect share at \textbf{36.2\,\%} (OC $\to$ DTTm $\to$ rhinoconjunctivitis),
with Co, Cu and Hg also partly routed through OP.
\textbf{Where it is weak:} cross-sectional symptom questionnaires, no effect sizes in the abstract, and OP assigned
from a sparse filter network --- a better fit may reflect smoother exposure error as much as biology.
\textbf{Why it matters here:} optical low-cost sensors report mass only; if OP is the health-relevant axis, mass
networks need a co-located speciation or OP tier to be interpretable for allergic endpoints.\par
\end{keybox}

\vspace{3mm}

\section*{\sffamily\bfseries\color{Deep}What changed today}
\vspace{-1.5mm}{\color{Amber}\rule{\linewidth}{1.1pt}}\vspace{2.5mm}

\footnotesize
\begin{itemize}

\item \textbf{A PM null built on interpolation.} Korean stroke registry, $n$=27,340: post-discharge district PM$_{10}$
and PM$_{2.5}$ (IDW) unrelated to recurrence; only SO$_2$ showed a graded association (Q4 vs Q1 HR \textbf{1.19},
0.94--1.50), concentrated in harbour cities.

\item \textbf{Secondary inorganic components and cognition.} CLHLS: per 10\,$\mu$g\,m$^{-3}$ HR 1.005 for PM$_{2.5}$
but 1.030 (sulphate) and 1.033 (ammonium); a favourable lifestyle halved risk and mediated $\sim$9\,\% of the
pollutant association.

\item \textbf{Traffic counts do not substitute for monitors.} Bosphorus bridges, 17,520\,h: a traffic-only proxy gave
negative R$^2$ for every pollutant; ML with meteorology and autoregressive terms reached R$^2\approx$0.77 ---
mostly persistence.

\item \textbf{Component-resolved assimilation.} CALIOP--MODIS species extinction assimilated by LETKF in WRF-Chem
recovers elevated free-tropospheric BC plumes over South Asia that the free model misses.

\item \textbf{One record held.} \emph{Foehn}, a generative km-scale global PM$_{2.5}$ forecaster (arXiv 2610.25941,
tier A), was admitted but its DataCite DOI does not yet resolve; it is authored and held for the next issue.

\end{itemize}

\vspace{2mm}

\begin{keybox}
{\sffamily\bfseries\footnotesize\color{Deep}Window continuity}\\[1.2mm]
{\fontsize{7.2}{8.8}\selectfont
The 18 October issue closed with \texttt{last\_entry\_date} at 18 October. This issue collects
\textbf{19 October in full}, continuous with it. Built on the run that started 23:0x CDT on 19 October (04:0x UTC
20 October), i.e.\ after the 22:00 cut-off, so the day is complete. 20 October is due on the next run at or after
22:00 CDT.\par}
\end{keybox}

\vspace{3mm}
\par\vskip 0pt plus 18\baselineskip\penalty-200\vskip 0pt plus -18\baselineskip
\section*{\sffamily\bfseries\color{Deep}Corpus analytics}
\vspace{-1.5mm}{\color{Amber}\rule{\linewidth}{1.1pt}}\vspace{2.5mm}

\noindent\begin{minipage}{\linewidth}\includegraphics[width=\linewidth]{\FIGDIR/f1_subtopics.png}
\figcap{Subtopic assignment of the 18 in-scope records across nine bands. Exposure assessment (5) and
burden/mitigation (4) lead; every health band holds a single record; other clinical endpoints is empty.}\end{minipage}

\vspace{2mm}
\noindent\begin{minipage}{\linewidth}\includegraphics[width=\linewidth]{\FIGDIR/f2_design_endpoint.png}
\figcap{Left --- architecture: two cohorts, two cross-sectional, three modelling/inventory, four
measurement, one laboratory; \textbf{metadata-only 6 (33\,\%)}, all Elsevier/T\&F titles without an indexed abstract
on day 0. Right --- endpoints: cognition, stroke recurrence, allergy symptoms, lung cancer proteome, urinary metals.}\end{minipage}

\vspace{2mm}
\noindent\begin{minipage}{\linewidth}\includegraphics[width=\linewidth]{\FIGDIR/f3_metric_geography.png}
\figcap{Left --- particle metric: unspeciated PM/emissions and PM$_{2.5}$-only records dominate;
one composition, one optical-property and one size-resolved record. Right --- geography: China 4, not-stated or
global 6, then Korea, East Asia, Iran, Turkiye, Thailand, Italy and the USA.}\end{minipage}

\vspace{2mm}
\noindent\begin{minipage}{\linewidth}\includegraphics[width=\linewidth]{\FIGDIR/f6_lifecourse.png}
\figcap{Life-course windows: every window populated by a single study except older adults (CLHLS
cognition, CPS-II smokers). The childhood and adolescent counts are the same Guangdong survey.}\end{minipage}

\vspace{2mm}
\noindent\begin{minipage}{\linewidth}\includegraphics[width=\linewidth]{\FIGDIR/f4_heatmap.png}
\figcap{Subtopic $\times$ study architecture. The only cell above one is exposure assessment $\times$
measurement campaign; the metadata-only column spreads across sensing, burden and indoor bands.}\end{minipage}

\vspace{2mm}

\noindent\begin{minipage}{\linewidth}\includegraphics[width=\linewidth]{\FIGDIR/f5_forest.png}
\figcap{Ratio scale: four HRs for incident cognitive disorder per 10\,$\mu$g\,m$^{-3}$ from one Chinese cohort
(PM$_{2.5}$, sulphate, nitrate, ammonium; single-pollutant models of collinear secondary components), and the
quartile HR for post-discharge SO$_2$ and stroke recurrence from a Korean registry, whose CI crosses 1. No PM
estimate from the stroke study is plotted because none was reported.}\end{minipage}

\vspace{4mm}

\par\vskip 0pt plus 24\baselineskip\penalty-200\vskip 0pt plus -24\baselineskip
\section*{\sffamily\bfseries\color{Deep}Paper-level digest}
\vspace{-1.5mm}{\color{Amber}\rule{\linewidth}{1.1pt}}\vspace{2.5mm}

{\fontsize{7.4}{9}\selectfont\color{Slate}Ordered by cluster. Each entry gives the
methodological core and the finding that matters, not an abstract paraphrase. Records
marked \textit{metadata only} carry no recoverable abstract and assert no finding.\par}

\band{Neuro / mental health\hfill 1 record}{Deep}

\paperentry{Zhang CLHLS PM$_{2.5}$ components cognitive disorders lifestyle mediation 2026}
{Association between long-term ambient air pollution exposure and cognitive disorders among older adults in China: the mediating role of lifestyle}
{J Glob Health}
{Chinese Longitudinal Healthy Longevity Study; Cox, Bayesian WQS and mediation models. Per 10 $\mu$g\,m$^{-3}$, PM$_{2.5}$ HR 1.005 (1.001-1.009), sulphate 1.030 (1.003-1.056), nitrate 1.020 (1.003-1.037), ammonium 1.033 (1.005-1.062), O$_3$ 1.006; mixture OR 1.028 (1.006-1.040) with O$_3$ carrying the largest WQS weight. Favourable vs unfavourable lifestyle HR 0.488 (0.336-0.710); lifestyle mediated 8.4-9.5\% of the pollutant associations. Limitation: an HR of 1.005 per 10 $\mu$g\,m$^{-3}$ is at the floor of detectability and the abstract gives neither n, case count nor how cognitive disorder was ascertained (MMSE cut-offs in CLHLS are education-sensitive); 'mediation' by lifestyle that is itself a confounder of residential exposure is not identifiable. Relevance: low - gridded component surfaces, no sensor component.}
{10.7189/jogh.16.04307}

\par\vskip 0pt plus 10\baselineskip\penalty-200\vskip 0pt plus -10\baselineskip
\band{Cardiovascular \& metabolic\hfill 1 record}{Teal}

\paperentry{Lee CRCS-K post-discharge SO$_2$ stroke recurrence PM null 2026}
{Post-Discharge Ambient Sulfur Dioxide Exposure and Stroke Recurrence After Acute Ischemic Stroke}
{J Stroke}
{27,340 acute ischaemic stroke patients (2014-2021), district-level IDW-interpolated pollutant means over the 3 months after discharge; cause-specific Cox with death as competing risk, cluster-robust SEs, splines. 764 recurrences, 542 deaths over a median 364.8 days. PM$_{10}$ and PM$_{2.5}$ showed no association; only SO$_2$ did (overall spline p=0.024), with Q4 vs Q1 HR 1.19 (0.94-1.50), stronger in age >=75 and women (interaction p 0.051, 0.062). Limitation: district IDW from sparse monitors gives Berkson-type error that biases PM toward the null, and SO$_2$ concentrates in harbour cities (Incheon, Ulsan, Busan) - a place effect cannot be excluded. Relevance: moderate - a null PM result built on interpolation is exactly where denser networks change the answer.}
{10.5853/jos.2026.03265}

\par\vskip 0pt plus 10\baselineskip\penalty-200\vskip 0pt plus -10\baselineskip
\band{Respiratory \& allergic\hfill 1 record}{Sky}

\paperentry{Wang Guangdong PM$_{2.5}$ oxidative potential DTT childhood allergy 2026}
{Associations of PM$_{2.5}$ and Its Oxidative Potential with Childhood Allergy-Related Symptoms: A Cross-Sectional Study in Guangdong Province, China}
{Environ Sci Technol}
{131,346 children aged 6-18 in six Guangdong cities. PM$_{2.5}$ mass and DTT oxidative potential were positively associated with all three symptom outcomes; mass-normalised OP (DTTm) gave the most consistent associations, lower AIC/BIC and the greatest stability in two-pollutant models. Mixture analyses implicated several constituents; OP plausibly carried part of the Co, OC, Cu and Hg associations, the largest indirect share 36.2\% for OC and rhinoconjunctivitis via DTTm. Limitation: cross-sectional symptom questionnaires; OP is presumably modelled or assigned from few filter sites, and the abstract reports no effect sizes, so the model-fit comparison cannot be weighed against exposure-error differences between metrics. Relevance: high - OP is not measurable by optical LCS, which bounds what a mass-only network can say about allergy.}
{10.1021/acs.est.6c05462}

\par\vskip 0pt plus 10\baselineskip\penalty-200\vskip 0pt plus -10\baselineskip
\band{Reproductive \& developmental\hfill 1 record}{Sage}

\paperentry{Moradnia Isfahan pregnant urinary metals indoor dust AQI 2026}
{Spatial variability of heavy metals in pregnant women: linking urinary biomarkers, indoor dust contamination, and air quality in Isfahan, Iran}
{Environ Geochem Health}
{80 pregnant women; paired urine and vacuum dust analysed by ICP-OES. Dust metals ranked Pb>Cr>Ni>Cd (Pb mean 111.10 ug/g); urine Pb 10.52, Cr 8.30, Ni 4.14 ug/L. Urinary metals and AQI differed across urban zones (p<0.001); dust did not (Pb p=0.514). Building age correlated with dust Pb (rho 0.788). Dust and urinary profiles were uncorrelated. Limitation: n=80, spot urine without creatinine adjustment stated, and the only airborne metric is a zone-level AQI - no PM mass or PM-bound metal was measured, so 'air quality' enters only as an ecological covariate. Relevance: low - illustrates the gap that personal or indoor PM monitoring with filter speciation would fill.}
{10.1007/s10653-026-03508-x}

\par\vskip 0pt plus 10\baselineskip\penalty-200\vskip 0pt plus -10\baselineskip
\band{Mechanistic toxicology\hfill 1 record}{Violet}

\paperentry{Tang CPS-II plasma proteome air pollution lung cancer ever-smokers 2026}
{Analysis of the Potential Mediating Role of Plasma Proteins in the Association Between Air Pollution Exposure and Lung Cancer Risk Among Ever-Smokers}
{Environ Sci Technol}
{95 matched case-control pairs of ever-smokers with pre-diagnostic plasma proteomics from the CPS-II Nutrition Cohort. Seventeen proteins associated with both exposures and lung cancer (CX3CL1, ADAM15, CXCL17 among them, p<0.05); high-dimensional mediation (HIMA) flagged six, including CX3CL1 and CXCL17 (Bonferroni p<0.05). Limitation: 190 people against a high-dimensional proteome, and the mediation signals disappeared when cases diagnosed within a year of blood draw were excluded - consistent with reverse causation by subclinical tumour. The abstract does not name the pollutants' exposure model. Pathways (chemokine signalling, immune modulation) are hypothesis-generating only. Relevance: low - biomarker discovery, no sensing implication.}
{10.1021/acs.est.6c07407}

\par\vskip 0pt plus 10\baselineskip\penalty-200\vskip 0pt plus -10\baselineskip
\band{Exposure assessment \& modelling\hfill 5 records}{Clay}

\paperentry{Cheng LETKF WRF-Chem CALIOP component-resolved aerosol assimilation East Asia 2026}
{Constraining the vertical profiles of aerosol component and heating rate over east Asia through assimilating CALIOP observations}
{Atmos Chem Phys}
{Component-resolved LETKF assimilation in WRF-Chem of CALIOP-MODIS synergistic species-specific extinction profiles (dust, sea salt, BC, water-soluble), with spatial and observational constraints. The analysis reconstructs multi-layer vertical structure and corrects severe underestimation of elevated BC plumes, including persistent free-tropospheric BC over South Asia; AERONET and AD-Net lidar validation shows gains in AOT, extinction and single-scattering albedo. Limitation: CALIOP's narrow curtain and 16-day repeat leave most of the domain unobserved, the species classification is retrieval-model dependent, and the abstract gives no quantitative skill figures; surface PM$_{2.5}$ is not evaluated. Relevance: moderate - column and vertical constraints are what satellite-surface fusion with LCS networks currently lacks.}
{10.5194/acp-26-14671-2026}

\paperentry{Xiao NCP Sichuan haze reduced-nitrogen NOC oligomerisation PM$_{2.5}$ 2026}
{Aqueous Carbonyl-Amino Aggregation Drives Atmospheric Nitrogen-Containing Oligomer Formation in China Haze}
{Environ Sci Technol}
{Winter-haze PM$_{2.5}$ from the North China Plain compared with the Sichuan Basin. Reduced-nitrogen NOCs were 40\% of ESI+ organics. High-molecular-weight (m/z>250) Re-NOCs rose in humid versus dry haze and correlated with NH4+ and the aqueous-SOA fraction of water-soluble organics, alongside longer aliphatic chains and higher unsaturation - read as carbonyl-amino oligomerisation and cyclisation. A marker oligomer (C8H13N3O2) tracked WSON; stronger acidity and higher temperature favoured oligomerisation over cyclisation, enriching it in the Sichuan Basin. Limitation: ESI+ is semi-quantitative and ionisation-biased, number of samples is not stated, and mechanism is inferred from correlation. Relevance: low - brown-carbon NOCs add absorbing mass that optical scattering sensors misweight under humid haze.}
{10.1021/acs.est.6c09940}

\paperentry{Lu street-view greenery street canyon PM concentrations 2026}
{Association between urban greenery and particulate matter concentrations in street canyons: a local-scale study based on street view imagery}
{Environ Sci: Atmos}
{Street-view imagery quantified street greenery quantity and configuration, which were related to PM concentrations in street canyons; the authors attribute the associations to low-level vegetation barriers restricting ventilation. The record's abstract is two sentences: no city, sample size, PM metric, measurement method or effect direction and magnitude are given. Limitation: on what is available, the inference is correlational and the proposed ventilation mechanism is untested; treat as near-metadata until the full text is read. Relevance: moderate - siting of street-level sensors relative to hedges and low canopy is a known representativeness issue for LCS networks, and image-derived greenery is a cheap covariate for spatial calibration.}
{10.1039/d6ea00086j}

\paperentry{He NYC-METS size-resolved organic aerosol offline AMS 2026}
{Characteristics of size-resolved organic aerosols using offline aerosol mass spectrometry: NYC-METS}
{Aerosol Sci Technol}
{\textsc{metadata only.} Title and authors only; no abstract in Crossref, Europe PMC or Semantic Scholar at harvest. The record reports size-resolved organic aerosol characterised by offline aerosol mass spectrometry - re-aerosolised filter or impactor extracts - within the NYC-METS campaign in New York City. Sampling period, size cuts, source apportionment and sample counts cannot be stated. Limitation: offline AMS loses water-insoluble and semi-volatile fractions in extraction, so size-resolved OA mass is a lower bound; no inference beyond the title is possible. Relevance: moderate - size-resolved organic composition governs refractive index and hygroscopicity, the two terms optical LCS corrections approximate.}
{10.1080/02786826.2026.2746516}

\paperentry{Barnett-Itzhaki hourly compound heat air pollution extremes 2026}
{Hourly Resolution Reveals Hidden Temporal Windows of Compound Heat-Air Pollution Extremes Masked by Daily Averaging}
{Atmos Pollut Res}
{\textsc{metadata only.} Title and author only; no abstract retrievable at harvest. The title claims that hourly data expose compound heat-air-pollution extreme windows that daily means hide. Pollutants, location, thresholds, record length and how many events are missed by daily averaging cannot be stated. Limitation: unverifiable beyond the title; compound-extreme definitions are threshold-dependent, and hourly co-occurrence of heat and O$_3$ is expected by photochemistry, so the PM-specific contribution is unknown. Relevance: moderate - sub-daily exposure windows are precisely what continuous low-cost monitors resolve and daily regulatory filters do not.}
{10.1016/j.apr.2026.103243}

\par\vskip 0pt plus 10\baselineskip\penalty-200\vskip 0pt plus -10\baselineskip
\band{Sensing, forecasting \& instrumentation\hfill 2 records}{Amber}

\paperentry{Birinci Bosphorus bridge traffic ML proxy PM$_{2.5}$ NO$_2$ 2026}
{Machine-Learning Proxy of Air-Quality Variability Associated with Intercontinental Bridge Traffic over the Bosphorus in Istanbul}
{Atmosphere}
{17,520-h panel (2022-2023) pairing six bridge traffic sensors with the nearest monitors. A Prophet traffic-only baseline gave negative R2 for every bridge x pollutant. LightGBM, XGBoost and a BiLSTM-attention model with meteorology, calendar, precipitation lags and short autoregressive terms reached mean hold-out R2 \textasciitilde{}0.77 (0.58-0.86), RMSE down 47-63\%; the three families were practically equivalent. SHAP ranked autocorrelation, season, wind direction and heavy vehicles above total traffic. Limitation: autoregressive pollutant terms make this nowcasting persistence, not a traffic proxy - the headline R2 largely reflects lagged concentrations, and 'closest monitor' distance is unstated. Relevance: moderate - shows traffic counts alone do not substitute for measurement in complex coastal flow.}
{10.3390/atmos17101024}

\paperentry{Ren physics-constrained UQ urban PM$_{2.5}$ prediction driving factors 2026}
{Integrating Implicit Physics Constraints and Uncertainty Quantification for Urban PM$_{2.5}$ Prediction and Driving Factor Analysis}
{Build Environ}
{\textsc{metadata only.} Title and authors only; no abstract retrievable at harvest. The title describes an urban PM$_{2.5}$ prediction model with implicit physics constraints and uncertainty quantification, plus driving-factor attribution. Architecture, city, input features, horizon, skill metrics and whether the uncertainty is calibrated cannot be stated. Limitation: 'implicit' physics constraints in ML usually mean loss-function regularisation, which does not guarantee physical consistency out of distribution; nothing beyond the title is verifiable. Relevance: moderate - calibrated predictive uncertainty is what allows LCS-fed forecasts to propagate sensor error honestly.}
{10.1016/j.buildenv.2026.115449}

\par\vskip 0pt plus 10\baselineskip\penalty-200\vskip 0pt plus -10\baselineskip
\band{Occupational \& indoor\hfill 2 records}{Amber}

\paperentry{Suh residential occupant activities airborne microplastic emissions 2026}
{Microplastics emissions from occupant activities in residential spaces: Experimental quantification and relative importance analysis of particle generation factors}
{Build Environ}
{\textsc{metadata only.} Title and authors only; no abstract retrievable at harvest. The title reports experimental quantification of microplastic particle emissions from occupant activities in residential spaces and a relative-importance analysis of generation factors (activity type, textiles, surfaces presumably). Size range, identification method (FTIR, Raman, pyrolysis GC-MS), emission rates and whether the inhalable fraction was measured cannot be stated. Limitation: unverifiable beyond the title; microplastic counts are method-bound and seldom expressed as PM mass. Relevance: moderate - indoor activity-resolved particle sources are what household LCS see as unexplained spikes; polymer fibres also scatter differently from combustion aerosol.}
{10.1016/j.buildenv.2026.115399}

\paperentry{Du classroom CO$_2$ proxy PM VOC ventilation types 2026}
{CO$_2$ as a proxy for PM and VOCs in classrooms across ventilation types and activities}
{Indoor Environments}
{\textsc{metadata only.} Title and authors only; no abstract retrievable at harvest. The title tests whether CO$_2$ tracks PM and VOCs in classrooms across ventilation types and activities - i.e. whether a cheap NDIR sensor can stand in for particle and VOC monitoring. Number of classrooms, ventilation categories, PM instrumentation and agreement statistics cannot be stated. Limitation: unverifiable beyond the title; physically, CO$_2$ tracks occupant-generated bioeffluents and air exchange, whereas PM indoors also depends on outdoor infiltration, resuspension and filtration, so decoupling under mechanical ventilation is expected. Relevance: high - directly asks whether school CO$_2$ monitor rollouts can be read as PM exposure surveillance.}
{10.1016/j.indenv.2026.100200}

\par\vskip 0pt plus 10\baselineskip\penalty-200\vskip 0pt plus -10\baselineskip
\band{Burden, policy \& mitigation\hfill 4 records}{Coral}

\paperentry{Zhang Jiaxing inventory co-control screening PM$_{2.5}$ GHG 2026}
{Inventory-Based Co-Control Screening of Air Pollutants and Greenhouse Gases in a Typical Yangtze River Delta City}
{Atmosphere}
{Harmonised 2024 integrated inventory plus enterprise statistics for Jiaxing; 54 non-zero source categories scored on an air-pollutant H-score (mean of PM$_{2.5}$, NOx, SO$_2$, VOCs, NH$_3$ shares) against a GWP100 CO$_2$-eq G-score, screened at the 70th percentile. Coupled co-control candidates: CHP, thermal power, motor vehicles; decoupled single-pollutant priorities: industrial solvents (VOCs) and paved-road dust (PM). A 75th-percentile benchmark implies 7,941 t VOC reduction among reporting enterprises. Limitation: equal-weight averaging of pollutant shares ignores toxicity and secondary PM formation potential; no ambient or health validation. Relevance: low - road-dust PM priority is the coarse fraction optical LCS undercount.}
{10.3390/atmos17101023}

\paperentry{Warminski SRC pellets boiler ESP PM CO emissions 2026}
{Short-Rotation Coppice Pellets as Biofuels in a Modern Boiler with an Electrostatic Precipitator: Particulate Matter and Carbon Monoxide Emissions}
{Molecules}
{Short-rotation coppice pellets (willow, poplar, black locust) versus pine and birch in a pellet boiler with an electrostatic precipitator. Species significantly affected PM and CO; the ESP cut mean PM by 52.5\% (53.81 to 25.54 mg/m3) without affecting CO. Highest PM: willow without ESP (90.96 mg/m3); lowest: pine with ESP (17.51). Poplar gave low PM but the highest CO; black locust the best balance (combustion efficiency 93.20\%). Limitation: one boiler, total PM only (no size distribution or PM$_{1}$ fraction, where residential-combustion mass sits), and steady-state operation that misses start-up peaks. Relevance: low - source-side mitigation; a 52\% ESP removal at total-mass level likely overstates sub-micron removal.}
{10.3390/molecules31193443}

\paperentry{Brunori Rome urban vineyard grapevine leaf PM interception 2026}
{Plant-Atmosphere Interactions in Urban Vineyards of Rome: A Digital Phenotyping Approach for the Functional Selection of Grapevine Cultivars}
{Plants (Basel)}
{Urban vineyard of 170 cultivars at Rome's Botanical Garden, 2025 season: 54 cultivars screened by leaf micromorphology, five assessed in depth (anatomy, chlorophyll, reflectance indices, SEM-counted PM). After \textasciitilde{}20 days without heavy rain, leaf-surface PM ranged from 52,024 (Barbera) to 77,462 particles/mm2 (Maturano), comparable to some urban tree species; adaxial deposition was over twice abaxial. Limitation: particle counts, not mass or size-resolved deposition flux; five cultivars, one site, one accumulation window, and no ambient PM measurement to convert counts into a removal rate. Relevance: low - green-infrastructure mitigation claims need co-located ambient PM, which a low-cost sensor at canopy height would supply cheaply.}
{10.3390/plants15193014}

\paperentry{Polwiang Bangkok Chiang Mai burning control DLNM GAM PM$_{2.5}$ regimes 2026}
{Differential sensitivity of urban and biomass-dominated PM$_{2.5}$ regimes to burning control: A distributed-lag generalized additive model analysis of Bangkok and Chiang Mai (2021-2024)}
{Atmos Pollut Res}
{\textsc{metadata only.} Title and author only; no abstract retrievable at harvest. The title states a distributed-lag GAM comparing the sensitivity of PM$_{2.5}$ to burning-control measures in an urban regime (Bangkok) and a biomass-burning-dominated regime (Chiang Mai), 2021-2024. Effect sizes, lag structure, the burning-control variable (fire counts, policy dates) and confounder set cannot be stated. Limitation: observational policy attribution in a GAM is sensitive to meteorological adjustment and transboundary smoke from Myanmar and Laos; single author. Relevance: moderate - Chiang Mai is a dense PMS-sensor setting (DustBoy network), and regime-specific response is where LCS spatial coverage tests policy effects.}
{10.1016/j.apr.2026.103236}

\clearpage
\section*{\sffamily\bfseries\color{Deep}Corpus register}
\vspace{-1.5mm}{\color{Amber}\rule{\linewidth}{1.1pt}}\vspace{2.5mm}

\input{table_rows}

\vspace{2mm}

\begin{keybox}
{\sffamily\bfseries\footnotesize\color{Deep}Provenance and method}\\[1.2mm]
{\fontsize{7.2}{8.8}\selectfont
Entry window \textbf{19 October 2026}, single day (Monday). Harvester legs: \textbf{PubMed} health
\textbf{4}, sensing \textbf{1}; \textbf{Europe PMC} \textbf{96} (16 AGRICOLA retro-index dropped); \textbf{Crossref by
ISSN} \textbf{55}; \textbf{arXiv} 1 entry dated 19 October; \textbf{OpenAlex} HTTP 429. 137 unique, 50 already
carried, 87 fresh. \textbf{PubMed connector} (EDAT 19 Oct) 6 PMIDs, 2 appended --- \textbf{89} candidates screened,
\textbf{19 admitted}, \textbf{70 rejected} with logged reasons; 1 admitted record \textbf{held} (arXiv DOI not yet
registered), leaving \textbf{18}. \textbf{Consensus} calibration sweep: 10 hits, 0 new. \textbf{DOI gate} 0 fail /
0 warn after the hold. \textbf{Trial watch}: 0 registrations updated 19--20 October. Abstract retry (PubMed, Europe
PMC, Crossref, Semantic Scholar) recovered none of the six day-0 Elsevier and T\&F records; \textbf{6 carried on
metadata alone}. \textbf{Scholar Gateway} not called. Abstracts remain the copyright of their publishers.\par}
\end{keybox}

\end{document}
